% 来源：OEBPS/Text/Section0001_split_005.xhtml → image_720e3dcce7e94eb1bdbddf2e961f5136.jpg
\begin{longtable}{@{}C{\dimexpr 0.50000\linewidth-1.00000\tabcolsep\relax}C{\dimexpr 0.50000\linewidth-1.00000\tabcolsep\relax}@{}}
\caption{标准普尔指数对比全天候投资组合（1939—2013年）}\label{tab:27-1}\\
\toprule\noalign{\vskip 3pt}
{\heiti 标准普尔指数} & {\heiti 全天候投资组合} \\
\midrule
\endfirsthead
\multicolumn{2}{r}{续表}\\[0.6em]
\toprule\noalign{\vskip 3pt}
{\heiti 标准普尔指数} & {\heiti 全天候投资组合} \\
\midrule
\endhead
\bottomrule
\endfoot
\bottomrule
\endlastfoot
75年里标准普尔指数亏损年度18个\footnote{\label{fn:table-表27-1标准普尔指数对比全天候投资组合-1}包括红利再投资} & 同期全天候投资组合亏损年度10个\footnote{\label{fn:table-表27-1标准普尔指数对比全天候投资组合-2}这10个亏损年度中有两年亏损的概率仅为0.03\%（基本上是收支平衡年，因此从实际操作的角度来看，75年里实际上只有8年是亏损的）}（平均略多于10年一次） \\
最大年度亏损43.3\% & 最大年度亏损3.93\% \\
平均年度亏损11.40\% & 平均年度亏损1.63\% \\
\end{longtable}
